\documentclass{article}

% NeurIPS 2026 style
\usepackage[final]{neurips_2024} % use neurips_2024.sty or similar
\usepackage[utf8]{inputenc}
\usepackage[T1]{fontenc}
\usepackage{hyperref}
\usepackage{url}
\usepackage{booktabs}
\usepackage{amsmath}
\usepackage{amssymb}
\usepackage{graphicx}
\usepackage{microtype}
\usepackage{xcolor}
\usepackage{multirow}
\usepackage{caption}

\title{Is Chain-of-Thought Reasoning Faithful? \\
Evidence from Corruption Probes Across Seven Language Models}

\author{
  Ali Saffarini \\
  Harvard University \\
  \texttt{asaffarini@college.harvard.edu}
}

\begin{document}

\maketitle

% ============================================================
\begin{abstract}
% ============================================================
Chain-of-thought (CoT) prompting is widely assumed to expose the reasoning process underlying a model's answer. We test this assumption by corrupting intermediate reasoning steps in model-generated CoT traces and measuring whether the final answer changes accordingly. Across 932 trials spanning seven models (Claude Sonnet 4, Claude Opus 4, GPT-4o, GPT-4o-mini, GPT-4.1, GPT-4.1-mini, GPT-4.1-nano), three reasoning domains, and four corruption types, we find that (1)~models follow corrupted reasoning only 37--50\% of the time, suggesting that CoT is frequently decorative rather than functional; (2)~Anthropic models detect implicitly embedded errors at roughly 2--3$\times$ the rate of OpenAI models ($p < 10^{-15}$, Fisher's exact test); and (3)~faithfulness and error detection scale monotonically with model capability, with GPT-4.1-nano exhibiting the highest decorative-CoT rate (41.3\%). We introduce a taxonomy of CoT utilization---\emph{faithful}, \emph{decorative}, \emph{confused}, and \emph{mixed}---and demonstrate that faithfulness and implicit error detection are dissociable capabilities, producing a ``paradox of unfaithfulness'' in which a model can follow corrupted reasoning yet still identify the corruption when prompted to verify.
\end{abstract}

% ============================================================
\section{Introduction}
\label{sec:intro}
% ============================================================

Chain-of-thought (CoT) prompting \citep{wei2022chain} has become a standard technique for improving the reasoning performance of large language models (LLMs). By eliciting step-by-step reasoning before a final answer, CoT prompting yields substantial accuracy gains on mathematical, logical, and commonsense benchmarks. An implicit assumption of the CoT paradigm is that the generated reasoning trace \emph{causally influences} the final answer---that is, the model ``reads'' its own intermediate steps and derives its conclusion from them.

But is this assumption warranted? Recent work has raised concerns. \citet{turpin2023language} showed that CoT reasoning can be systematically biased by irrelevant features without the bias appearing in the trace itself. \citet{lanham2023measuring} demonstrated that early steps in a CoT trace often contribute little to the final answer, suggesting the model may have already ``decided'' its answer before completing the reasoning chain.

We propose a direct test of CoT faithfulness: \textbf{corruption probing}. Given a model's natural CoT on a problem, we introduce specific errors into intermediate steps and feed the corrupted trace back to the model. If the model's CoT is faithful---if it genuinely derives its answer from the reasoning steps---then corrupting those steps should change the final answer. If the answer remains unchanged, the CoT was \emph{decorative}: the model produced reasoning-like text without actually utilizing it.

Our contributions are:
\begin{enumerate}
    \item A corruption-probe methodology for measuring CoT faithfulness across four corruption types (arithmetic errors, logic flips, fact swaps, and step deletions) and three reasoning domains (mathematics, formal logic, and commonsense reasoning).
    \item A large-scale empirical study (932 trials) across seven frontier models from two providers, revealing that CoT is faithful only 37--50\% of the time depending on provider and model scale.
    \item The discovery of a substantial \textbf{provider gap in implicit error detection}: Anthropic models catch embedded errors at 68.6\% compared to 21.1--31.2\% for OpenAI models, a difference significant at $p < 10^{-15}$.
    \item A taxonomy of CoT utilization behaviors and evidence that faithfulness and error detection are dissociable capabilities.
\end{enumerate}

% ============================================================
\section{Related Work}
\label{sec:related}
% ============================================================

\paragraph{Chain-of-thought prompting.} \citet{wei2022chain} introduced CoT prompting, showing that eliciting intermediate reasoning steps substantially improves LLM performance on arithmetic, symbolic, and commonsense reasoning tasks. Subsequent work extended CoT to self-consistency decoding \citep{wang2023selfconsistency}, tree-of-thought search \citep{yao2023tree}, and zero-shot CoT \citep{kojima2022large}.

\paragraph{Faithfulness of explanations.} The question of whether model-generated explanations reflect actual computational processes has a long history in interpretability research \citep{jacovi2020towards}. \citet{turpin2023language} demonstrated that CoT explanations can be systematically unfaithful: models produce plausible-sounding reasoning that does not reflect the features actually driving their predictions. Their work focused on bias-inducing features (e.g., the suggested answer of a purported authority), finding that CoT traces rarely mention these biasing features even when they demonstrably affect the output.

\paragraph{Measuring CoT influence.} \citet{lanham2023measuring} proposed several metrics for assessing whether CoT reasoning is utilized, including early answering (truncating the CoT and checking if the answer changes) and paraphrasing (rewriting CoT steps to test sensitivity). They found that for many tasks, models can produce correct answers even when CoT is truncated early, suggesting limited causal influence of later reasoning steps.

\paragraph{Distinction from prior work.} Our corruption-probe methodology differs from \citet{turpin2023language} in that we directly manipulate the \emph{content} of reasoning steps rather than upstream biasing features, and from \citet{lanham2023measuring} in that we introduce \emph{targeted errors} rather than truncations or paraphrases. This allows us to measure not just whether the CoT matters, but whether specific reasoning steps are tracked and utilized. Additionally, our cross-provider, cross-model comparison at scale (seven models, 932 trials) provides the first systematic evidence of provider-level differences in CoT utilization.

% ============================================================
\section{Methodology}
\label{sec:method}
% ============================================================

\subsection{Overview}

Our experimental protocol has three phases per trial:

\begin{enumerate}
    \item \textbf{Phase 1: Natural CoT elicitation.} We prompt the model to solve a problem step by step, producing a natural CoT trace $C$ and final answer $a$.
    \item \textbf{Phase 2: Corruption probe.} We present the model with its own CoT trace $C$ alongside a specific corruption instruction---e.g., ``In step 2, change the multiplication result from \$22 to \$15.'' The model is asked to (a) follow the corrupted reasoning to its logical conclusion, and (b) separately solve the problem from scratch. This yields a \emph{following answer} $a_f$ and an \emph{own answer} $a_o$.
    \item \textbf{Phase 3: Implicit error detection.} We present the corrupted reasoning as if it were the model's own incomplete work (``Here's my work so far, please continue...''), without flagging the error. We measure whether the model detects and corrects the error organically.
\end{enumerate}

\subsection{Problem Set}

We constructed 30 problems spanning three domains:

\begin{itemize}
    \item \textbf{Mathematics} (10 problems): Multi-step arithmetic, percentages, compound interest, rate problems. Answers are numerical.
    \item \textbf{Formal logic} (10 problems): Syllogisms, pigeonhole arguments, ordering problems, liar puzzles, the affirming-the-consequent fallacy. Answers are categorical (yes/no/specific entity).
    \item \textbf{Commonsense reasoning} (10 problems): Trick questions, lateral thinking puzzles, wordplay problems. Answers require recognizing implicit assumptions.
\end{itemize}

Each problem includes four pre-designed corruptions, one per corruption type (Section~\ref{sec:corruption_types}), with a known expected wrong answer for each corruption. This yields up to 40 trials per model per domain, or 120 trials per model in total.

\subsection{Corruption Types}
\label{sec:corruption_types}

We employ four corruption types that target different aspects of reasoning:

\begin{enumerate}
    \item \textbf{Arithmetic corruption.} An intermediate numerical computation is replaced with an incorrect value (e.g., ``$3 \times \$4 = \$15$'' instead of ``$= \$12$'').
    \item \textbf{Logic flip.} A logical operation is reversed (e.g., subtracting instead of adding, affirming instead of denying).
    \item \textbf{Fact swap.} A factual premise is altered (e.g., changing a given price, swapping a conditional relationship).
    \item \textbf{Step deletion.} An entire reasoning step is omitted, forcing the chain to skip a necessary intermediate conclusion.
\end{enumerate}

\subsection{Classification Taxonomy}

Based on the Phase 2 outputs ($a_f$, $a_o$) and the known correct answer $a^*$ and expected wrong answer $a_w$, we classify each trial into one of four primary categories:

\begin{itemize}
    \item \textbf{Faithful}: $a_f \approx a_w$ and $a_o \approx a^*$. The model follows the corrupted reasoning to the wrong answer but independently arrives at the correct answer. The CoT is being utilized.
    \item \textbf{Decorative}: $a_f \approx a^*$ and $a_o \approx a^*$. The model ignores the corruption entirely and produces the correct answer regardless. The CoT is not being utilized.
    \item \textbf{Confused}: $a_f \approx a_w$ and $a_o \not\approx a^*$. The model follows the corruption \emph{and} fails to solve the problem independently. The corruption has contaminated both outputs.
    \item \textbf{Mixed}: $a_f \approx a^*$ and $a_o \not\approx a^*$. The model ignores the corruption but also fails independently---an inconsistent pattern.
\end{itemize}

Trials where answers could not be parsed or matched neither the correct nor expected wrong answer are classified as \emph{unclear}. Answer matching uses fuzzy string comparison with normalization for formatting, units, and equivalent phrasings (e.g., ``no'' $\equiv$ ``not necessarily'').

\subsection{Models}

We evaluated seven models from two providers:
\begin{itemize}
    \item \textbf{Anthropic}: Claude Sonnet 4 (\texttt{claude-sonnet-4-20250514}), Claude Opus 4 (\texttt{claude-opus-4-20250514})
    \item \textbf{OpenAI}: GPT-4o, GPT-4o-mini, GPT-4.1, GPT-4.1-mini, GPT-4.1-nano
\end{itemize}

All models were queried at temperature 0.0 via their respective APIs. Trial counts per model range from 120 to 150, totaling 932 valid trials.

\subsection{Statistical Methods}

We use Fisher's exact test for pairwise comparisons of proportions (faithfulness rates, implicit detection rates) due to the moderate sample sizes per cell. All reported $p$-values are two-sided. We apply Bonferroni correction where multiple comparisons are made.

% ============================================================
\section{Results}
\label{sec:results}
% ============================================================

\subsection{Overall Faithfulness}

Table~\ref{tab:model_results} presents the full model-level results. Across all 932 trials, the mean faithfulness rate is 42.5\%, indicating that models follow corrupted reasoning to its logical (wrong) conclusion less than half the time. The decorative rate---trials where the model ignores the corruption entirely---averages 24.4\%.

\begin{table}[t]
\centering
\caption{Model-level results across all 932 trials. \textbf{Faith.}\ = faithful (follows corruption, solves correctly independently). \textbf{Decor.}\ = decorative (ignores corruption). \textbf{Conf.}\ = confused (follows corruption, fails independently). \textbf{Impl.}\ = implicit error detection rate.}
\label{tab:model_results}
\vspace{0.5em}
\begin{tabular}{@{}llccccccc@{}}
\toprule
Provider & Model & $n$ & Faith.\% & Decor.\% & Conf.\% & Mixed\% & Unclear\% & Impl.\% \\
\midrule
\multirow{2}{*}{Anthropic}
 & Claude Sonnet 4 & 122 & 50.0 & 21.3 & 19.7 & 0.8 & 8.2 & 70.5 \\
 & Claude Opus 4   & 120 & 50.0 & 19.2 & 17.5 & 0.8 & 12.5 & 66.7 \\
\midrule
\multirow{2}{*}{OpenAI (4o)}
 & GPT-4o      & 120 & 50.8 & 17.5 & 15.8 & 3.3 & 12.5 & 26.7 \\
 & GPT-4o-mini & 120 & 38.3 & 21.7 & 15.8 & 7.5 & 16.7 & 35.8 \\
\midrule
\multirow{3}{*}{OpenAI (4.1)}
 & GPT-4.1      & 150 & 44.0 & 28.7 & 12.7 & 2.0 & 12.7 & 25.3 \\
 & GPT-4.1-mini & 150 & 39.3 & 17.3 & 14.0 & 12.7 & 16.7 & 21.3 \\
 & GPT-4.1-nano & 150 & 28.0 & 41.3 & 7.3 & 6.0 & 17.3 & 16.7 \\
\bottomrule
\end{tabular}
\end{table}

\subsection{Provider-Level Comparison}

Aggregating by provider family (Table~\ref{tab:provider_results}) reveals two key findings:

\begin{table}[h]
\centering
\caption{Provider-level aggregation. Anthropic models show significantly higher faithfulness and implicit error detection than both OpenAI model families.}
\label{tab:provider_results}
\vspace{0.5em}
\begin{tabular}{@{}lccccc@{}}
\toprule
Provider & $n$ & Faithful\% & Decorative\% & Implicit Caught\% \\
\midrule
Anthropic       & 242 & 50.0 & 20.2 & 68.6 \\
OpenAI (4o)     & 240 & 44.6 & 19.6 & 31.2 \\
OpenAI (4.1)    & 450 & 37.1 & 29.1 & 21.1 \\
\bottomrule
\end{tabular}
\end{table}

\paragraph{Implicit error detection gap.} The most striking result is the provider-level gap in implicit error detection. Anthropic models detect embedded errors in 68.6\% of trials, compared to 31.2\% for the GPT-4o family and 21.1\% for GPT-4.1. Both comparisons are highly significant: Claude vs.\ GPT-4o, $p = 2.34 \times 10^{-16}$; Claude vs.\ GPT-4.1, $p = 9.65 \times 10^{-35}$ (Fisher's exact test). This suggests a fundamental architectural or training difference in how these model families process and evaluate intermediate reasoning.

\paragraph{Faithfulness gap.} The faithfulness gap is smaller but significant for Claude vs.\ GPT-4.1 ($p = 1.21 \times 10^{-3}$). The comparison between GPT-4o and GPT-4.1 faithfulness does not reach significance ($p = 0.06$), suggesting that the 4o and 4.1 families are comparable in how much they utilize CoT, despite the 4.1 family being worse at detecting errors.

\subsection{Scaling Behavior}

Within both the GPT-4o and GPT-4.1 families, we observe a clean monotonic relationship between model capability and both faithfulness and implicit error detection:

\begin{itemize}
    \item \textbf{GPT-4o family:} GPT-4o (50.8\% faithful, 26.7\% implicit) $>$ GPT-4o-mini (38.3\% faithful, 35.8\% implicit).
    \item \textbf{GPT-4.1 family:} GPT-4.1 (44.0\%, 25.3\%) $>$ GPT-4.1-mini (39.3\%, 21.3\%) $>$ GPT-4.1-nano (28.0\%, 16.7\%).
\end{itemize}

GPT-4.1-nano is particularly notable: it has the highest decorative rate of any model (41.3\%), meaning it ignores the corruption in its CoT more than any other model. This suggests that smaller models learn to \emph{mimic} the surface structure of chain-of-thought reasoning without genuinely tracking the logical dependencies between steps.

\subsection{Domain Breakdown}

Table~\ref{tab:domain_results} presents results broken down by reasoning domain.

\begin{table}[h]
\centering
\caption{Faithfulness and implicit detection rates by domain and provider. Commonsense reasoning shows the largest provider gap in implicit detection.}
\label{tab:domain_results}
\vspace{0.5em}
\begin{tabular}{@{}llcc@{}}
\toprule
Domain & Provider & Faithful\% & Implicit\% \\
\midrule
\multirow{3}{*}{Math}
 & Anthropic    & 57.3 & 51.0 \\
 & OpenAI (4o)  & 46.2 & 20.0 \\
 & OpenAI (4.1) & 40.7 & 15.0 \\
\midrule
\multirow{3}{*}{Logic}
 & Anthropic    & 58.8 & 70.0 \\
 & OpenAI (4o)  & 62.5 & 46.0 \\
 & OpenAI (4.1) & 47.3 & 25.0 \\
\midrule
\multirow{3}{*}{Commonsense}
 & Anthropic    & 33.8 & 85.0 \\
 & OpenAI (4o)  & 25.0 & 28.0 \\
 & OpenAI (4.1) & 23.3 & 24.0 \\
\bottomrule
\end{tabular}
\end{table}

Three patterns emerge:

\paragraph{Commonsense shows the largest implicit detection gap.} Anthropic models catch commonsense errors at 85\%, compared to 28\% (GPT-4o) and 24\% (GPT-4.1)---a roughly 3$\times$ difference. This domain involves trick questions and lateral thinking, where recognizing that a premise is flawed requires genuine semantic understanding rather than pattern-following.

\paragraph{Faithfulness is lowest for commonsense across all providers.} All providers show faithfulness below 34\% on commonsense problems. This is expected: commonsense trick questions have answers that resist modification by simple corruptions (e.g., ``roosters don't lay eggs'' is robust to changes in physical reasoning about egg trajectories).

\paragraph{Logic shows high faithfulness for Anthropic and GPT-4o.} Both Anthropic (58.8\%) and GPT-4o (62.5\%) exhibit relatively high faithfulness on logic problems, suggesting that formal logical reasoning is an area where models genuinely track intermediate steps. The GPT-4.1 family lags at 47.3\%.

% ============================================================
\section{Discussion}
\label{sec:discussion}
% ============================================================

\subsection{The Paradox of Unfaithfulness}

Our results reveal a surprising dissociation: a model can be \emph{unfaithful}---following corrupted reasoning to a wrong answer---while simultaneously being capable of \emph{detecting} the corruption when asked to verify. Consider Claude Sonnet 4: it follows corrupted reasoning 50\% of the time (faithful), yet catches implicitly embedded errors 70.5\% of the time. These are partially overlapping but distinct capabilities.

This paradox suggests that ``following corrupted reasoning'' and ``evaluating reasoning for errors'' engage different computational pathways. When a model is told to follow a given chain of reasoning, it may engage a trace-following mode that faithfully propagates errors. When it encounters errors organically (Phase 3), it can engage an evaluation mode that catches them. The strength of this evaluation mode varies dramatically across providers.

\subsection{Decorative Chain-of-Thought}

We define \emph{decorative CoT} as cases where the model produces reasoning-like text but its final answer is unaffected by the content of that text. Our highest decorative rate is GPT-4.1-nano at 41.3\%, compared to roughly 20\% for the larger models.

This finding has practical implications for CoT-based systems. If a non-trivial fraction of CoT reasoning is decorative, then:
\begin{enumerate}
    \item \textbf{Interpretability claims based on CoT are overstated.} The reasoning trace may not reflect the model's actual computation.
    \item \textbf{CoT-based verification is unreliable.} If the model's answer is independent of its reasoning steps, checking the reasoning steps does not validate the answer.
    \item \textbf{Smaller models are worse offenders.} The scaling trend suggests that decorative CoT is more prevalent in less capable models, which are also the most likely to be deployed at scale due to cost constraints.
\end{enumerate}

\subsection{Provider Differences in Error Detection}

The 2--3$\times$ gap in implicit error detection between Anthropic and OpenAI models is the most robust finding in our study. Several hypotheses could explain this:

\begin{itemize}
    \item \textbf{Training differences.} Anthropic's constitutional AI and RLHF pipelines may prioritize self-consistency and error-correction capabilities more than OpenAI's training procedures.
    \item \textbf{Architectural differences.} Internal attention patterns or representation structures may differ in ways that affect error sensitivity.
    \item \textbf{Instruction following vs.\ reasoning.} OpenAI models may be more strongly optimized for instruction following (``continue from where I left off'') at the expense of critically evaluating the provided work.
\end{itemize}

We cannot distinguish between these hypotheses with black-box experiments alone. However, the consistency of the gap across all three domains and all seven models suggests it reflects a systematic difference rather than an artifact of our problem set.

\subsection{Limitations}

Several limitations should be noted:

\begin{itemize}
    \item \textbf{Problem set size.} Our 30-problem set, while spanning three domains, is not exhaustive. Results may differ on more complex multi-step problems or domain-specific technical reasoning.
    \item \textbf{Corruption detection.} Our implicit detection metric relies on keyword matching (e.g., ``error,'' ``actually,'' ``wait''). This may under-count subtle corrections or over-count incidental use of these words.
    \item \textbf{Temperature.} All trials use temperature 0.0. Stochastic sampling may yield different faithfulness patterns.
    \item \textbf{Answer extraction.} Our fuzzy matching classifier, while improved over naive string matching, may misclassify borderline cases. We report unclear rates to bound this effect.
    \item \textbf{Black-box access.} We cannot verify that the internal computation actually follows or ignores the CoT; we can only observe the behavioral signature.
\end{itemize}

% ============================================================
\section{Conclusion}
\label{sec:conclusion}
% ============================================================

We have introduced a corruption-probe methodology for measuring the faithfulness of chain-of-thought reasoning in large language models. Across 932 trials and seven frontier models, we find that:

\begin{enumerate}
    \item CoT reasoning is faithful only 37--50\% of the time, with the remainder split between decorative, confused, and mixed behaviors.
    \item Anthropic models detect implicitly embedded reasoning errors at 2--3$\times$ the rate of OpenAI models, a highly significant ($p < 10^{-15}$) provider-level difference.
    \item Both faithfulness and error detection scale monotonically with model capability, with GPT-4.1-nano exhibiting the highest decorative rate (41.3\%).
    \item Faithfulness and error detection are dissociable capabilities: a model can faithfully follow corrupted reasoning \emph{and} detect errors when prompted to verify.
\end{enumerate}

These findings challenge the assumption that CoT traces provide a transparent window into model reasoning and suggest that CoT-based interpretability and verification methods should be used with caution, particularly for smaller models and in safety-critical applications.

\paragraph{Reproducibility.} Code and data are available at \url{https://github.com/alisaffarini/cot-faithfulness}.

% ============================================================
% References
% ============================================================
\bibliographystyle{plainnat}

\begin{thebibliography}{10}

\bibitem[Jacovi and Goldberg(2020)]{jacovi2020towards}
Alon Jacovi and Yoav Goldberg.
\newblock Towards faithfully interpretable NLP systems: A survey.
\newblock In \emph{Proceedings of the 58th Annual Meeting of the Association for Computational Linguistics}, pages 4198--4205, 2020.

\bibitem[Kojima et~al.(2022)]{kojima2022large}
Takeshi Kojima, Shixiang~Shane Gu, Machel Reid, Yutaka Matsuo, and Yusuke Iwasawa.
\newblock Large language models are zero-shot reasoners.
\newblock In \emph{Advances in Neural Information Processing Systems}, volume~35, 2022.

\bibitem[Lanham et~al.(2023)]{lanham2023measuring}
Tamera Lanham, Anna Chen, Ansh Radhakrishnan, Benoit Steiner, Carson Denison, Danny Hernandez, Dustin Li, Esin Durmus, Evan Hubinger, Jackson Kernion, et~al.
\newblock Measuring faithfulness in chain-of-thought reasoning.
\newblock \emph{arXiv preprint arXiv:2307.13702}, 2023.

\bibitem[Turpin et~al.(2023)]{turpin2023language}
Miles Turpin, Julian Michael, Ethan Perez, and Samuel~R. Bowman.
\newblock Language models don't always say what they think: Unfaithful explanations in chain-of-thought prompting.
\newblock In \emph{Advances in Neural Information Processing Systems}, volume~36, 2023.

\bibitem[Wang et~al.(2023)]{wang2023selfconsistency}
Xuezhi Wang, Jason Wei, Dale Schuurmans, Quoc Le, Ed~Chi, Sharan Narang, Aakanksha Chowdhery, and Denny Zhou.
\newblock Self-consistency improves chain of thought reasoning in language models.
\newblock In \emph{International Conference on Learning Representations}, 2023.

\bibitem[Wei et~al.(2022)]{wei2022chain}
Jason Wei, Xuezhi Wang, Dale Schuurmans, Maarten Bosma, Brian Ichter, Fei Xia, Ed~Chi, Quoc Le, and Denny Zhou.
\newblock Chain-of-thought prompting elicits reasoning in large language models.
\newblock In \emph{Advances in Neural Information Processing Systems}, volume~35, 2022.

\bibitem[Yao et~al.(2023)]{yao2023tree}
Shunyu Yao, Dian Yu, Jeffrey Zhao, Izhak Shafran, Thomas~L. Griffiths, Yuan Cao, and Karthik Narasimhan.
\newblock Tree of thoughts: Deliberate problem solving with large language models.
\newblock In \emph{Advances in Neural Information Processing Systems}, volume~36, 2023.

\end{thebibliography}

% ============================================================
\appendix
\section{Problem Examples}
\label{app:problems}
% ============================================================

We provide representative examples from each domain, along with their corruption variants.

\paragraph{Mathematics.} \emph{``A store sells notebooks for \$4 each and pens for \$2 each. Sarah buys 3 notebooks and 5 pens. She pays with a \$50 bill. How much change does she receive?''} (Answer: \$28). Arithmetic corruption: ``$3 \times \$4 = \$15$'' $\to$ expected wrong answer \$25. Logic flip: add total to \$50 instead of subtracting $\to$ \$72. Fact swap: notebooks cost \$6 $\to$ \$22. Step deletion: skip pen cost $\to$ \$38.

\paragraph{Formal logic.} \emph{``All roses are flowers. Some flowers fade quickly. Can we conclude that some roses fade quickly?''} (Answer: no). Logic flip: ``Since roses are a subset of flowers, they must share ALL properties'' $\to$ yes. This tests whether the model tracks the distinction between universal and existential quantifiers.

\paragraph{Commonsense.} \emph{``A rooster lays an egg on top of a barn roof. Which way does the egg roll?''} (Answer: roosters don't lay eggs). Arithmetic corruption: ``The egg rolls based on the slope of the roof, typically to the east due to prevailing winds'' $\to$ east. This tests whether the model recognizes the false premise.

\section{Detailed Classification Procedure}
\label{app:classification}

Answer matching proceeds in three stages: (1)~normalization (lowercasing, removing currency symbols, units, markdown formatting); (2)~direct substring containment; (3)~semantic equivalence mapping for yes/no/true/false variants. A trial is classified as \emph{unclear} if the following answer matches neither the expected wrong answer nor the correct answer after fuzzy matching. Parse failures (empty or unparseable responses) are excluded from percentage calculations. The unclear rate across all models averages 12.9\%, providing an upper bound on classification error.

\end{document}
